% LaTeX version of alpha-winding-2026-09-24.typ, for journals and arXiv.
% The Typst file is the reference text; keep the two in step when either changes.
% Build from this directory with pdflatex -output-directory=<a scratch dir> run three times, so no build files land here.
% The checks are research/verify_alpha_winding.py, with output in research/artifacts/verify_alpha_winding.txt.
% The figure is research/figures/alpha-winding.pdf, written by plot_alpha_winding.py.
\documentclass[11pt]{article}
\usepackage[letterpaper,margin=1in]{geometry}
\usepackage[T1]{fontenc}
\usepackage{lmodern}
\usepackage{amsmath,amssymb,amsthm}
\usepackage{graphicx}
\usepackage{array}
\usepackage{booktabs}
\usepackage{microtype}
\usepackage[hidelinks]{hyperref}

\hypersetup{
  pdftitle={A sharp winding bound for the self-similar collapse of three point vortices in the alpha-models},
  pdfauthor={Chase Hendrick}}

\newtheorem{theorem}{Theorem}
\newtheorem{lemma}{Lemma}
\newtheorem{corollary}{Corollary}

\newcommand{\re}{\operatorname{Re}}
\newcommand{\im}{\operatorname{Im}}
\newcommand{\C}{\mathbb{C}}

\title{\Large\bfseries A sharp winding bound for the self-similar collapse\\
  of three point vortices in the $\alpha$-models}
\author{Chase Hendrick\\[0.2em] {\small GENChase}}
\date{}

\begin{document}
\maketitle

\begin{abstract}
In the $\alpha$-models of two-dimensional flow a point vortex of circulation $\Gamma$ induces the velocity $\Gamma r^{-\alpha-1}/(2\pi)$ at distance $r$; $\alpha = 0$ is the Euler equation and $\alpha = 1$ the surface quasi-geostrophic equation. In a self-similar collapse of three point vortices every vortex moves on a logarithmic spiral about the collision point, and the rotation per unit decrease of the logarithm of the squared size, $P = |\im\kappa|/(2|\re\kappa|)$, measures how tightly the spiral winds. We prove that $P > \sqrt{3+\alpha}/(2+\alpha)$ for every self-similar collapse of three point vortices when $\alpha > -0.896$, and that the constant is sharp: it is approached, but not attained, as two weak vortices of opposite sign merge next to a strong one. Equivalently, every vortex travels more than $(4+\alpha)/(2+\alpha)$ times its initial distance from the collision point, and its path makes an angle larger than $\arccos((2+\alpha)/(4+\alpha))$ with the direction to that point. The bound decreases from $\sqrt{3}/2$ in the Euler case to $2/3$ in the surface quasi-geostrophic case and tends to $0$ as $\alpha \to \infty$. The proof rests on two elementary identities: the circulations of a collapsing triangle are fixed, up to a factor, by its side lengths, and $P$ is a symmetric function of the side lengths alone. For $\alpha = 0$ it gives a short proof of the Euler bound $\sqrt{3}/2$. For four or more vortices in the Euler case the bound fails: numerically, four vortices reach $P = 0.7978967838\ldots$ and sixty-one vortices $P = 0.4981844392\ldots$.

\medskip
\noindent\textbf{Keywords:} point vortices; vortex collapse; self-similar motion; surface quasi-geostrophic equation; generalized Euler equation.\\
\textbf{MSC 2020:} 76B47, 76U60, 37N10.
\end{abstract}

\section{Introduction}

Three point vortices whose angular impulse and harmonic combination of circulations vanish can shrink to a point in finite time while the triangle they form keeps its shape and rotates. Each vortex then moves on a logarithmic spiral about the collision point. In the Euler case this self-similar collapse goes back to Gr\"obli; Conte and de Seze \cite{cds1980} and Kimura \cite{kimura1987} wrote the motion in closed form, and the companion preprint \cite{hendrick2026} shows that the winding of the spiral, measured by the dimensionless number $P$ of \eqref{eq:P} below, satisfies $P > \sqrt{3}/2$ for every self-similar collapse of three vortices, with a sharp constant.

Self-similar collapse of three vortices also occurs in the family of active-scalar models in which the streamfunction is obtained from the vorticity by a fractional inverse Laplacian, the generalized Euler or $\alpha$-models, which include the surface quasi-geostrophic (SQG) equation \cite{bb2018,gp2025,reinaud2021,rds2022}. Their point vortices move by
\begin{equation}\label{eq:bs}
\frac{dz_j}{dt} = \frac{i}{2\pi} \sum_{k\ne j} \Gamma_k \frac{z_j - z_k}{|z_j - z_k|^{2\beta}}, \qquad \beta = 1 + \frac{\alpha}{2},
\end{equation}
with positions $z_j \in \C$ and real circulations $\Gamma_j$. The case $\alpha = 0$ ($\beta = 1$) is the Euler system, and $\alpha = 1$ ($\beta = 3/2$) is SQG. A constant factor in front of the sum, which differs between papers, only rescales time and does not change $P$. In the conventions of Badin and Barry \cite{bb2018} and of Grotto and Pappalettera \cite{gp2025} the same system has the parameter $2 - \alpha$.

This note extends the winding bound to these models. With $B(\alpha) = \sqrt{3+\alpha}/(2+\alpha)$ we prove:

\begin{theorem}\label{thm:main}
Let $\alpha > \alpha_0$, where $\alpha_0 = -0.8960874362\ldots$ is given by \eqref{eq:alpha0}. Every self-similar collapse of three point vortices under \eqref{eq:bs} has $P > B(\alpha)$, and for every $\delta > 0$ there is one with $P < B(\alpha) + \delta$.
\end{theorem}

The proof (Section~\ref{sec:bound}) needs two facts proved in Section~\ref{sec:rates}: a collapsing triangle is scalene and carries circulations fixed up to a factor by its side lengths (Lemma~\ref{lem:circ}), and $P$ is a symmetric function of the side lengths alone (Lemma~\ref{lem:rates}). For $\alpha = 0$ it gives a short proof of the Euler bound of \cite{hendrick2026}, though not the minimum for each ratio of circulations obtained there. Section~\ref{sec:more} records that the bound does not extend to more than three vortices in the Euler case, and Section~\ref{sec:numerics} lists the numerical checks.

\section{Self-similar motion}

Call a motion of \eqref{eq:bs} \emph{self-similar} if the triangle keeps its shape and orientation, that is, $z_j(t) = a(t) + b(t) z_j(0)$ with complex $a(t)$ and $b(t) \ne 0$, and \emph{collapsing} if moreover $b(t) \to 0$ in finite time.

\begin{lemma}\label{lem:ss}
Let a self-similar motion of \eqref{eq:bs} change size. Then $\Gamma_1 + \Gamma_2 + \Gamma_3 \ne 0$, and with $z_c = \sum_j \Gamma_j z_j/\sum_j \Gamma_j$ the center of vorticity, there is $\kappa_0 \in \C$ with $\re\kappa_0 \ne 0$ such that $\dot z_j = \kappa(t)(z_j - z_c)$ with $\kappa(t) = \kappa_0 \lambda^{-2\beta}$, and
\begin{equation}\label{eq:spiral}
z_j(t) - z_c = \lambda(t) e^{i\phi(t)} (z_j(0) - z_c), \qquad \lambda^{2\beta} = 1 + 2\beta \re\kappa_0\, t, \qquad \phi = \frac{\im\kappa_0}{2\re\kappa_0} \ln\lambda^2 .
\end{equation}
The motion collapses if and only if $\re\kappa_0 < 0$, at $t_c = -1/(2\beta\re\kappa_0)$. Every vortex moves on a logarithmic spiral about $z_c$; the number
\begin{equation}\label{eq:P}
P = \frac{|\im\kappa|}{2|\re\kappa|}
\end{equation}
does not depend on $t$, and it is the angle through which the triangle turns while its squared size decreases by the factor $e$. The path of each vortex makes the angle $\arctan 2P$ with the direction to $z_c$, a vortex at distance $r_0$ from $z_c$ travels the distance $r_0\sqrt{1 + 4P^2}$ before the collapse, and the initial angular velocity $\omega_0 = \im\kappa_0$ satisfies $|\omega_0| t_c = P/\beta$.
\end{lemma}

\begin{proof}
The linear impulse $Q = \sum_j \Gamma_j z_j$ is conserved. If $\sum_j \Gamma_j = 0$, then $Q(t) = b(t) Q(0)$, and since $|b|$ changes, $Q(0) = 0$; then $\Gamma_1(z_1 - z_3) + \Gamma_2(z_2 - z_3) = 0$ with real coefficients not both zero, so the vortices are collinear, and a collinear triangle does not change size under \eqref{eq:bs} (Lemma~\ref{lem:circ} below). Hence $\sum_j \Gamma_j \ne 0$, $z_c$ is conserved, and $z_c = a + b z_c$, so $z_j(t) - z_c = b(t) w_j$ with $w_j = z_j(0) - z_c$. The right side of \eqref{eq:bs} is homogeneous: replacing $z_j - z_c$ by $b w_j$ multiplies it by $b|b|^{-2\beta}$. Writing $\dot z_j(0) = \kappa_0 w_j$, which holds for the same $\kappa_0$ for all $j$ because $\dot b(0) w_j = \dot z_j(0)$, we get $\dot b = \kappa_0 b |b|^{-2\beta}$. With $b = \lambda e^{i\phi}$ this reads $\lambda^{2\beta - 1}\dot\lambda = \re\kappa_0$ and $\dot\phi = \im\kappa_0\,\lambda^{-2\beta}$, which integrate to \eqref{eq:spiral}. The size changes, so $\re\kappa_0 \ne 0$. Along the spiral $r = r_0\lambda$ and $d\phi/d\ln\lambda = 2P$ in absolute value, so the tangent of the angle between the path and the radius is $2P$ and the arc length from $r_0$ to $0$ is $r_0\sqrt{1 + 4P^2}$. Finally $|\omega_0| t_c = |\im\kappa_0|/(2\beta|\re\kappa_0|) = P/\beta$.
\end{proof}

The spiral form \eqref{eq:spiral} of a self-similar collapse in the $\alpha$-models is in \cite{yoi2021} and \cite[Eqs.~(3)--(6)]{iy2023}, in \cite[(A.19)]{dgc2023}, whose H\"older exponent $1/(\alpha + 1)$ is our $1/(2\beta)$, and in \cite[Prop.~2.1]{gp2025}. For $\beta = 1$, $P = |\omega_0| t_c$ is the product studied in \cite{hendrick2026}. For $\beta \ne 1$ the product $|\omega_0| t_c$ is the rotation per unit decrease of $\ln(t_c - t)$, and it differs from $P$ by the factor $1/\beta$; we state the bound for $P$, which measures the geometry of the spiral, and for $|\omega_0| t_c$ in Corollary~\ref{cor:geom}.

\section{Circulations and rates of a collapsing triangle}\label{sec:rates}

Label the sides by the opposite vertices: $r_1 = |z_2 - z_3|$, $r_2 = |z_3 - z_1|$, $r_3 = |z_1 - z_2|$, and put $f_i = r_i^{-2\beta}$. Let $A_s = \im(\overline{(z_2 - z_1)}(z_3 - z_1))/2$ be the signed area, positive when $z_1, z_2, z_3$ are counterclockwise, and $A = |A_s|$. Indices $(i, j, k)$ run over the cyclic permutations of $(1, 2, 3)$.

\begin{lemma}\label{lem:circ}
Under \eqref{eq:bs}, $d(r_i^2)/dt = (2A_s/\pi)\,\Gamma_i (f_j - f_k)$. A motion is self-similar and changes size if and only if the triangle is not collinear, its sides are pairwise different, and
\begin{equation}\label{eq:circ}
\Gamma_i = c\, \frac{r_i^2}{f_j - f_k} \qquad (i = 1, 2, 3)
\end{equation}
for some real $c \ne 0$. In that case $\re\kappa = cA_s/\pi$, and exactly one of a triangle and its mirror image collapses.
\end{lemma}

\begin{proof}
For $i = 3$, the terms of $\dot z_1 - \dot z_2$ proportional to $z_1 - z_2$ do not change $|z_1 - z_2|$, and the rest is $(i\Gamma_3/2\pi)((z_1 - z_3) f_2 - (z_2 - z_3) f_1)$. Since $\im(\overline{(z_1 - z_2)}(z_1 - z_3)) = \im(\overline{(z_1 - z_2)}(z_2 - z_3)) = 2A_s$, $d(r_3^2)/dt = 2\re(\overline{(z_1 - z_2)}(\dot z_1 - \dot z_2)) = (2A_s/\pi)\Gamma_3(f_1 - f_2)$, and the other sides follow by cyclic relabeling. In a self-similar motion $d\ln r_i^2/dt = 2\re\kappa$ for every $i$, so $(2A_s/\pi)\Gamma_i(f_j - f_k)/r_i^2$ is the same for all $i$; if the size changes, this common value is not $0$, so $A_s \ne 0$, $f_j \ne f_k$ for all pairs, and \eqref{eq:circ} holds with $c = \pi\re\kappa/A_s$. Conversely, if \eqref{eq:circ} holds with $c \ne 0$ on a scalene triangle that is not collinear, the three logarithmic rates coincide, so the ratios of the squared sides are stationary at this shape; the rates of these ratios are functions of the shape alone, multiplied by a power of the size, so the shape is an equilibrium of the reduced motion and is kept for all time, with the orientation, since the area cannot pass through zero. The motion is therefore self-similar, and its size changes because $\re\kappa = cA_s/\pi \ne 0$. Reflecting the triangle changes the sign of $A_s$ and keeps \eqref{eq:circ}, so it changes the sign of $\re\kappa$.
\end{proof}

In particular, isosceles, equilateral and collinear triangles never collapse, and every scalene triangle collapses in exactly one orientation, with the circulations \eqref{eq:circ} fixed up to a real factor. On a scalene triangle, \eqref{eq:circ} says that the vector $(1/\Gamma_i)$ is proportional to the cross product of $(r_i^2)$ and $(r_i^{2 - 2\beta})$, that is, orthogonal to both; so it is equivalent to the vanishing of the angular impulse, $\sum_i r_i^2/\Gamma_i = 0$, together with $\sum_i r_i^{-\alpha}/\Gamma_i = 0$, which is the vanishing of $\sum_{j<k}\Gamma_j\Gamma_k r_i^{-\alpha}$ for $\alpha \ne 0$ and the classical condition $\sum_{j<k}\Gamma_j\Gamma_k = 0$ for $\alpha = 0$. Conditions of this kind are given in \cite{bb2018,cl2024,gp2025,reinaud2021,rds2022}; Lemma~\ref{lem:circ} is a short self-contained form of them, with the circulations explicit.

\begin{lemma}\label{lem:rates}
For a self-similar motion that changes size, with \eqref{eq:circ},
\begin{equation}\label{eq:S}
2\pi\im\kappa = \frac12 \sum_{i=1}^3 \Gamma_i (f_j + f_k) = \frac{c}{2} S, \qquad S = \sum_{i=1}^3 r_i^2 \coth\Bigl(\beta\ln\frac{r_k}{r_j}\Bigr), \qquad P = \frac{|S|}{8A}.
\end{equation}
\end{lemma}

\begin{proof}
Since $\dot z_1 - \dot z_2 = \kappa(z_1 - z_2)$, $\im\kappa = \im((\dot z_1 - \dot z_2)/(z_1 - z_2))$. By \eqref{eq:bs},
\[
2\pi\frac{\dot z_1 - \dot z_2}{z_1 - z_2} = i(\Gamma_1 + \Gamma_2) f_3 + i\Gamma_3 \frac{f_2(z_1 - z_3) - f_1(z_2 - z_3)}{z_1 - z_2},
\]
and by the law of cosines $\re((z_1 - z_3)/(z_1 - z_2)) = (r_2^2 + r_3^2 - r_1^2)/(2r_3^2)$ and $\re((z_2 - z_3)/(z_1 - z_2)) = (r_2^2 - r_1^2 - r_3^2)/(2r_3^2)$. Hence
\[
2\pi\im\kappa = (\Gamma_1 + \Gamma_2) f_3 + \Gamma_3 \frac{f_1 + f_2}{2} + \Gamma_3 \frac{(f_2 - f_1)(r_2^2 - r_1^2)}{2r_3^2}.
\]
By \eqref{eq:circ} the last term is $c(r_1^2 - r_2^2)/2$, and also $c(r_1^2 - r_2^2) = \Gamma_1(f_2 - f_3) + \Gamma_2(f_1 - f_3)$. Substituting gives $2\pi\im\kappa = \frac12\sum_i \Gamma_i(f_j + f_k)$, and by \eqref{eq:circ}, $\Gamma_i(f_j + f_k) = c r_i^2 (f_j + f_k)/(f_j - f_k) = c r_i^2 \coth(\beta\ln(r_k/r_j))$. With $\re\kappa = cA_s/\pi$ from Lemma~\ref{lem:circ}, \eqref{eq:P} gives $P = |S|/(8A)$.
\end{proof}

The circulations drop out: $P$ depends only on the shape of the triangle. Interchanging two labels reverses the cyclic order, which changes the sign of every term of $S$, so $|S|$ is a symmetric function of the side lengths. The first identity says that the rotation rate is half the sum of the rates at which the three pairs would rotate in isolation.

\section{The bound}\label{sec:bound}

Put
\begin{equation}\label{eq:alpha0}
\beta_0 = \frac{3 + \sqrt{105}}{24} = 0.5519562819\ldots, \qquad \alpha_0 = 2\beta_0 - 2 = -0.8960874362\ldots,
\end{equation}
so that $\beta_0$ is the positive root of $12\beta^2 - 3\beta - 2$. We use the elementary inequality
\begin{equation}\label{eq:log}
\frac{1}{\ln(1 + x)} \ge \frac1x + \frac12 - \frac{x}{12} \qquad (x > 0).
\end{equation}
It is trivial where the right side is negative, that is, for $x \ge 3 + \sqrt{21}$. For $0 < x < 3 + \sqrt{21}$, the function $g(x) = 12x/(12 + 6x - x^2) - \ln(1 + x)$ has $g(0) = 0$ and $g'(x) = x^3(24 - x)/((12 + 6x - x^2)^2(1 + x)) > 0$, so $\ln(1 + x) \le 12x/(12 + 6x - x^2)$, which is \eqref{eq:log}.

\begin{proof}[Proof of Theorem~\ref{thm:main}]
By Lemma~\ref{lem:circ} a collapsing triangle is scalene, and by Lemma~\ref{lem:rates}, $P = |S|/(8A)$ is symmetric in the sides and invariant under scaling. Order and scale the sides as $r_2 = \rho < r_3 = 1 < r_1$, and write $r_1^2 = 1 + \rho m$. With $\psi$ the angle at $z_1$, between the sides of lengths $\rho$ and $1$, the law of cosines gives $2\cos\psi = \rho - m$, so $0 < m < 2 + \rho$, $0 < \rho < 1$, $A = (\rho/2)\sin\psi$ and $4\sin^2\psi = 4 - (m - \rho)^2$. Put $X = \rho^{2\beta}$, $Y = (1 + \rho m)^\beta$ and $b = (\beta/2)\ln(1 + \rho m)$. Then $0 < X < 1 < Y$ and, from \eqref{eq:S} with $(i, j, k) = (1, 2, 3)$ and its cyclic shifts,
\[
S = (1 + \rho m)\frac{1 + X}{1 - X} + \rho^2\coth b - \frac{Y + X}{Y - X}.
\]
(i) The function $t \mapsto (t + X)/(t - X)$ decreases for $t > X$, so $(Y + X)/(Y - X) \le (1 + X)/(1 - X)$ and $S \ge \rho m + \rho^2\coth b$. (ii) $\coth b > 1/b$, and \eqref{eq:log} with $x = \rho m$ gives $\rho^2/b \ge (\rho/\beta)(2/m + \rho - \rho^2 m/6)$. Hence
\[
S > \frac{\rho}{\beta}\Lambda, \qquad \Lambda = \beta m + \frac{2}{m} + \rho - \frac{\rho^2 m}{6} > 0,
\]
the last because $\rho^2 m/6 < \rho/2$. In particular $S > 0$, and
\[
P = \frac{S}{4\rho\sin\psi} > \frac{\Lambda}{2\beta\sqrt{4 - (m - \rho)^2}}.
\]
(iii) The identity
\begin{equation}\label{eq:id}
\Lambda^2 - (1 + 2\beta)(4 - (m - \rho)^2) = \Bigl((1 + \beta)m - \frac2m - \rho\Bigr)^2 + \frac{\rho^2}{3}\Bigl(1 + 6\beta - \beta m^2 - \rho m + \frac{\rho^2 m^2}{12}\Bigr)
\end{equation}
holds for all $\rho$, $m$, $\beta$, and its right side is positive. If $m < 2$, the second bracket exceeds $1 + 6\beta - 4\beta - 2 = 2\beta - 1 > 0$. If $m \ge 2$, then $(1 + \beta)m - 2/m - \rho \ge 1 + 2\beta - \rho > 2\beta$, so the square exceeds $4\beta^2$; and since $m < 2 + \rho$, the second term is at least $h(\rho)/3$ with $h(\rho) = \rho^2(1 + 6\beta - \beta(2 + \rho)^2 - \rho(2 + \rho))$, where $h(\rho) - h(1) = (1 - \rho)(\beta\rho^3 + 5\beta\rho^2 + 3\beta\rho + 3\beta + \rho^3 + 3\rho^2 + 2\rho + 2) > 0$ and $h(1) = -3\beta - 2$; so the right side exceeds $4\beta^2 - \beta - 2/3 = (12\beta^2 - 3\beta - 2)/3 > 0$ for $\beta > \beta_0$. Hence $\Lambda > \sqrt{1 + 2\beta}\sqrt{4 - (m - \rho)^2}$ and $P > \sqrt{1 + 2\beta}/(2\beta) = B(\alpha)$.

Sharpness. Let $\rho \to 0$ with $m$ fixed in $(0, 2)$. Since $2\beta > 1$, $X = o(\rho)$, so the first and last terms of $S$ are $1 + \rho m + o(\rho)$ and $1 + o(\rho)$, while $\rho^2\coth b = 2\rho/(\beta m) + O(\rho^2)$. Hence $S/\rho \to m + 2/(\beta m)$ and
\[
P \to P_0(m) = \frac{\beta m^2 + 2}{2\beta m\sqrt{4 - m^2}},
\]
whose only critical point on $(0, 2)$ is $m_* = \sqrt{2/(1 + \beta)}$, with $P_0(m_*) = \sqrt{1 + 2\beta}/(2\beta)$. By Lemma~\ref{lem:circ} each of these triangles, or its mirror image, collapses with the circulations \eqref{eq:circ}, so $P$ comes arbitrarily close to $B(\alpha)$.
\end{proof}

In the extremal limit $\rho \to 0$ the circulations \eqref{eq:circ} are, up to a factor, $\Gamma_2 \approx 1$ at the vertex opposite the shortest side and $\Gamma_1 \approx -\Gamma_3$ of order $\rho^{2\beta - 1}$ at its ends: two weak vortices of opposite sign, at distance $\rho$, spiral into a strong vortex, and the line through them makes the angle $\arccos(1/\sqrt{2(1 + \beta)})$ with the line to the strong vortex. For $\alpha = 0$ this is the limit $\mu \to 0$ of \cite[Corollary~1]{hendrick2026}.

\begin{corollary}\label{cor:geom}
Let $\alpha > \alpha_0$. In every self-similar collapse of three point vortices under \eqref{eq:bs}, each vortex travels more than $(4 + \alpha)/(2 + \alpha)$ times its initial distance from the collision point, the angle between its path and the direction to that point exceeds $\arccos((2 + \alpha)/(4 + \alpha))$, and $|\omega_0| t_c > 2\sqrt{3 + \alpha}/(2 + \alpha)^2$. None of these constants can be improved.
\end{corollary}

\begin{proof}
By Lemma~\ref{lem:ss} the three quantities are $\sqrt{1 + 4P^2}$, $\arctan 2P$ and $P/\beta$, which increase with $P$, and $\sqrt{1 + 4B^2} = (4 + \alpha)/(2 + \alpha)$.
\end{proof}

\begin{table}[ht]
\centering
\begin{tabular}{lccccc}
\toprule
 & $\alpha = 0$ (Euler) & $\alpha = 1/2$ & $\alpha = 1$ (SQG) & $\alpha = 3/2$ & $\alpha = 2$ \\
\midrule
$B(\alpha)$ & $\sqrt3/2 = 0.8660$ & $\sqrt{14}/5 = 0.7483$ & $2/3$ & $0.6061$ & $\sqrt5/4 = 0.5590$ \\
path ratio & $2$ & $9/5$ & $5/3$ & $11/7$ & $3/2$ \\
angle & $60^\circ$ & $56.25^\circ$ & $53.13^\circ$ & $50.48^\circ$ & $48.19^\circ$ \\
\bottomrule
\end{tabular}
\caption{The bound of Theorem~\ref{thm:main} and the two constants of Corollary~\ref{cor:geom} for some values of $\alpha$.}
\label{tab:values}
\end{table}

\begin{figure}[ht]
\centering
\includegraphics[width=\textwidth]{figures/alpha-winding.pdf}
\caption{(a) The bound $B(\alpha)$ of Theorem~\ref{thm:main} (curve) and $P = |S|/(8A)$ for 24395 random collapsing triangles at 61 values of $\alpha$ (dots, spread slightly in $\alpha$ for visibility); none lies below the curve, which is approached as the shortest side tends to $0$. (b) The paths of the three vortices of an SQG collapse ($\alpha = 1$) with sides $0.45$, $1$ and $(1 + 0.45 m_*)^{1/2}$ and the circulations of Lemma~\ref{lem:circ}, integrated from \eqref{eq:bs} to $10^{-6} t_c$ before the collapse; dots mark the initial positions and the cross the collision point. The program is \texttt{plot\_alpha\_winding.py}.}
\label{fig:alpha}
\end{figure}

Since $d\ln B/d\alpha = -(4 + \alpha)/(2(3 + \alpha)(2 + \alpha)) < 0$, the bound decreases strictly in $\alpha$, from $\sqrt3/2$ in the Euler case through $2/3$ for SQG to $0$ as $\alpha \to \infty$: more local interactions allow tighter spirals, but never a collapse without rotation (Figure~\ref{fig:alpha}). The restriction $\alpha > \alpha_0$ comes only from the case $m \ge 2$ in step (iii); numerical minimization of \eqref{eq:S} for $\alpha$ between $-1.8$ and $\alpha_0$, in arithmetic whose precision grows with $|\ln\rho|$, finds no value below $B(\alpha)$ either, the infimum again being approached only as $\rho \to 0$, but we do not claim the bound there.

\section{More than three vortices}\label{sec:more}

Self-similar motion of any number of vortices has the same form \cite{dk2014,gotoda2021,oneil2007}, $\dot z_j = \kappa(z_j - z_c)$ and $P$ defined by \eqref{eq:P}. In the Euler case the bound of \cite{hendrick2026} does not extend to more vortices: the seven-vortex collapse of Demina and Kudryashov \cite[Table~1]{dk2014} has $P = 12433/(1240\sqrt{155}) = 0.805\ldots$ \cite{hendrick2026}. A numerical minimization of $\im\kappa/(-2\re\kappa)$ over the circulations and positions of four vortices, from more than two hundred random starts, ends at a single configuration up to relabeling, with
\[
P_4 = 0.79789678387986348076760587182\ldots,
\]
a strict local minimum (the Hessian of the Lagrangian on the three-dimensional tangent space of the constraints has eigenvalues $2.21$, $12.58$ and $18.75$). In the normalization $\Gamma_1 = 1$, $z_c = 0$, $\kappa = -1 + ib$ it has $b = 1.5957935677597269615\ldots$ and
\begin{align*}
\Gamma &= (1,\ 0.1673914318538,\ 0.1041331880362,\ -0.2272513003829),\\
z &= (0.0166495669172,\ -0.2191226402232 - 0.3290676220068\,i,\\
&\qquad -0.1579874300122 - 0.2864759934272\,i,\ -0.1605333843048 - 0.3736601672808\,i).
\end{align*}
Families of larger numerical minimizers go lower still: $P = 0.5172291322\ldots$ for 33 vortices and $P = 0.4981844392\ldots$ for 61 vortices, so a self-similar collapse can turn by less than half a radian per unit decrease of $\ln\lambda^2$. Whether $P$ can be arbitrarily small for large numbers of vortices is open; the families computed so far extrapolate to about $0.477$. These are numerical results, not proofs, and we do not know whether the three-vortex bound extends in some modified form to $\alpha \ne 0$ for more vortices.

\section{Numerical verification}\label{sec:numerics}

The program \texttt{verify\_alpha\_winding.py} checks each step in high-precision arithmetic, and its output is \texttt{artifacts/verify\_alpha\_winding.txt}. On 700 random scalene triangles with $\alpha \in \{0, 0.3, 1, 1.7, 3\}$, the circulations \eqref{eq:circ} make the Biot--Savart velocities \eqref{eq:bs} self-similar to $10^{-38}$, and $P$ computed from the velocities agrees with \eqref{eq:S} to $3\cdot 10^{-38}$. The derivative of $g$, the identity \eqref{eq:id} and the factorization of $h(\rho) - h(1)$ are verified symbolically, and the critical point and value of $P_0$ as well. The chain of inequalities in the proof holds at $20000$ random points with $\beta$ from $0.56$ to $20$ and $\rho$ down to $10^{-8}$, computed with precision adapted to $\rho^{2\beta}$; in double precision the first step appears to fail at small $\rho$ and large $\beta$, where its margin, of order $\rho^{2\beta}$, is below the rounding error. For $\alpha$ from $-1.8$ to $-0.85$ the minimum of $P - B(\alpha)$ over the ordered triangles is positive in every run, again only with precision adapted to $\rho$. For $\alpha = 1$, the scalene triangles on which \eqref{eq:circ} gives two equal circulations carry the ratios $-\Gamma_2/\Gamma_1$ from $0.3876$ to $0.49999$ on a grid, in agreement with the interval of \cite[Lemma~1]{bb2018}, and at $\Gamma = 0.49$ the ratio of sides is $0.7515$, against $0.751484$ in \cite{bb2018}. Near-extremal triangles at $\rho = 10^{-6}$ and $m = m_*$ collapse, and $P - B(\alpha)$ is $1.3\cdot 10^{-12}$, $1.8\cdot 10^{-13}$, $1.6\cdot 10^{-13}$ and $1.6\cdot 10^{-13}$ for $\alpha = 0, 1/2, 1, 2$. The four-, thirty-three- and sixty-one-vortex configurations of Section~\ref{sec:more} satisfy the self-similarity equations to $10^{-29}$, $10^{-39}$ and $10^{-28}$ relative in independent Biot--Savart checks.

\section{Discussion}

Self-similar collapse of three vortices in the $\alpha$-models has been studied from several directions: Badin and Barry \cite{bb2018} derive the necessary conditions for collapse in Nambu form and show that in the SQG case collapse is self-similar when the energy vanishes and not self-similar otherwise; for circulations $(1, -\Gamma, 1)$ they find self-similar SQG collapse exactly for $0.387464\ldots < \Gamma < 1/2$ \cite[Lemma~1]{bb2018}, an interval that Lemma~\ref{lem:circ} reproduces; Reinaud, Dritschel and Scott \cite{rds2022} give the conditions for self-similar collapse of three vortices and map the collapse time over the collapsing configurations; Reinaud \cite{reinaud2021} gives, for SQG, two necessary conditions and the collapse time in terms of the side lengths (as quoted in \cite{cl2024}), which has the form of the rate $\re\kappa$ of Lemma~\ref{lem:circ} at $\alpha = 1$; Chen and Liu \cite{cl2024} prove that the existence of self-similar collapse and expansion depends only on the circulations, give explicit expressions for all nontrivial self-similar solutions \cite[Theorems~2.1, 2.2]{cl2024}, and obtain the interval of Badin and Barry analytically \cite[Theorem~4.1]{cl2024}, a description from the side of the circulations that Lemma~\ref{lem:circ} complements from the side of the shape; Yasunaga, Otobe and Iwayama \cite{yoi2021} give necessary conditions and the orbits of self-similar motion, and Iwayama and Yajima \cite{iy2023} show that the collapsing motions are linearly unstable and the expanding ones stable; Iwayama, Yajima and Watanabe \cite{iyw2025} prove that collapse that is not self-similar can occur for $0 < \alpha < 2$ and cannot for $-1 \le \alpha < 0$; and Grotto and Pappalettera \cite[Prop.~2.1]{gp2025} characterize the self-similar motions and write the collapsing spiral, as a step toward bursts out of arbitrary configurations. We have not found a bound on the winding, the spiral angle or the distance traveled in these works, nor the explicit form \eqref{eq:circ} of the circulations or the formula \eqref{eq:S} for $P$, but we have read only the abstract and introduction of \cite{cl2024} and have not been able to consult \cite{reinaud2021}. In the Euler case the closed forms of the rates go back to \cite{cds1980} and \cite{kimura1987}, and the bound $P > \sqrt3/2$ with its minimum for each ratio of circulations is in \cite{hendrick2026}; Theorem~\ref{thm:main} gives, for $\alpha = 0$, a shorter proof of that bound through the side lengths.

The bound concerns self-similar collapse only; it says nothing about other singular motions of the $\alpha$-models, such as the collapses at nonzero energy found in \cite{bb2018,iyw2025}. Its value lies in the geometry: however the circulations are chosen, a collapsing triangle turns by more than $B(\alpha)$ radians per unit decrease of the logarithm of its squared size, and the infimum is realized only by a degenerate dipole beside a strong vortex.

\medskip
{\sloppy\noindent\textbf{Data availability.} The programs \texttt{verify\_alpha\_winding.py} and \texttt{plot\_alpha\_winding.py} and their output are in the directory \texttt{research/} of the repository \url{https://github.com/SharpMeow/GENChase}; the many-vortex configurations and their checks are in \texttt{research/generalizations-2026-09-24/}.\par}

\smallskip
\noindent{\small This work was prepared with AI assistance. The author takes full responsibility for its content.}

\begin{thebibliography}{15}

\bibitem{bb2018} G. Badin and A. M. Barry, Collapse of generalized Euler and surface quasigeostrophic point vortices, \emph{Phys. Rev. E} \textbf{98} (2018) 023110; arXiv:1805.10127.

\bibitem{cl2024} J. Chen and Q. Liu, Sufficient and necessary conditions for self-similar motions of three point vortices in generalized fluid systems, \emph{Physica D} \textbf{470} (2024) 134392.

\bibitem{cds1980} R. Conte and L. de Seze, Exact solution of the planar motion of three arbitrary point vortices, report DPhG/PSRM/1697/80, CEN Saclay (1980); arXiv:1511.00069.

\bibitem{dk2014} M. V. Demina and N. A. Kudryashov, Rotation, collapse, and scattering of point vortices, \emph{Theor. Comput. Fluid Dyn.} \textbf{28} (2014) 357--368; doi:10.1007/s00162-014-0319-4.

\bibitem{dgc2023} M. Donati and L. Godard-Cadillac, H\"older regularity for collapses of point vortices, \emph{Nonlinearity} \textbf{36} (2023) 5773; arXiv:2111.14230.

\bibitem{gotoda2021} T. Gotoda, Self-similar motions and related relative equilibria in the $N$-point vortex system, \emph{J. Dyn. Differ. Equ.} \textbf{33} (2021) 1759--1777; arXiv:2002.09624.

\bibitem{gp2025} F. Grotto and U. Pappalettera, Collapse and burst of generalized surface quasi-geostrophic point vortices, preprint, arXiv:2505.19782 (2025).

\bibitem{hendrick2026} C. Hendrick, Minimal winding in the self-similar collapse of three point vortices and of two concentric vortex polygons, preprint (2026), \url{https://github.com/SharpMeow/GENChase}.

\bibitem{iy2023} T. Iwayama and T. Yajima, Linear stability of self-similar motions of three point vortices in a generalized two-dimensional fluid system, \emph{J. Phys. Soc. Jpn.} \textbf{92} (2023) 084401.

\bibitem{iyw2025} T. Iwayama, S. Yajima and T. Watanabe, Non-self-similar collapse of three point vortices in a generalized two-dimensional fluid system, \emph{J. Phys. A} \textbf{58} (2025) 075701.

\bibitem{kimura1987} Y. Kimura, Similarity solution of two-dimensional point vortices, \emph{J. Phys. Soc. Jpn.} \textbf{56} (1987) 2024--2030.

\bibitem{oneil2007} K. A. O'Neil, Relative equilibrium and collapse configurations of four point vortices, \emph{Regul. Chaotic Dyn.} \textbf{12} (2007) 117--126.

\bibitem{reinaud2021} J. N. Reinaud, Self-similar collapse of three geophysical vortices, \emph{Geophys. Astrophys. Fluid Dyn.} \textbf{115} (2021) 369--392.

\bibitem{rds2022} J. N. Reinaud, D. G. Dritschel and R. K. Scott, Self-similar collapse of three vortices in the generalised Euler and quasi-geostrophic equations, \emph{Physica D} \textbf{434} (2022) 133226.

\bibitem{yoi2021} T. Yasunaga, N. Otobe and T. Iwayama, \emph{J. Phys. Soc. Jpn.} \textbf{90} (2021) 124401.

\end{thebibliography}

\end{document}
